% F087｜牛顿切线法快速求近似根｜01_statement.tex
% 继承上传模板的 LaTeX 环境；当前条目的内容独立编写。
\clearpage
\begin{statementbox}
\textbf{\textcolor{CStatement}{为什么用切线找根}}
\par\medskip
\begin{center}
% F087｜tangent 原创简明教学示意图
\begin{tikzpicture}[>=stealth,scale=.82]
\draw[->] (0,0)--(7.6,0) node[right]{$x$};\draw[->] (.8,-1.4)--(.8,3.2) node[above]{$y$};
% 横向比例尺 x_point=8(x-2.65)，纵向比例尺 y_point=0.75(x^2-10)。
\draw[teal,thick,domain=2.76:3.48,samples=45] plot ({8*(\x-2.65)},{.75*(\x*\x-10)});
% 在真实点 x_0=3 的切线；交轴位置 x_1=19/6。
\draw[orange,thick] (2,-1.2)--(6,1.05);
\fill[orange] (2.8,-.75) circle (2pt) node[below left=5pt]{$x_0=3$};
\fill[orange] (4.1333,0) circle(2pt) node[below=7pt]{$x_1\approx3.1667$};
\node at (5.5,2.7){\small 用切线寻找新的交点};
\node[teal] at (5.7,1.9){\small $y=x^2-10$};
\end{tikzpicture}

\end{center}
\medskip
方程 $x^2-10=0$ 的正根在 $3$ 与 $4$ 之间。把曲线在 $x=3$ 附近用一小段切线代替，切线与横轴的交点就给出新的猜测值。重复做，通常能越来越接近真正的根。
\end{statementbox}
\clearpage
\begin{statementbox}
\textbf{\textcolor{CStatement}{三个最有用的公式}}
\par\medskip
\begin{center}
% F087｜square 原创简明教学示意图
\begin{tikzpicture}[>=stealth,scale=.9]
\node[draw,fill=teal!8,rounded corners,minimum width=7.3cm,minimum height=1.5cm,align=center] at (4,0) {$x_{n+1}=\dfrac12\left(x_n+\dfrac{A}{x_n}\right)$};
\node at (4,-1.2) {\small 正数 $x_n$ 与 $A/x_n$ 的平均数};
\end{tikzpicture}

\end{center}
\medskip
\[x_{n+1}=x_n-\frac{f(x_n)}{f'(x_n)}.\]求平方根：$x_{n+1}=\dfrac12(x_n+A/x_n)$；求正数的立方根：$x_{n+1}=\dfrac{2x_n+A/x_n^2}{3}$。后两个只是第一个公式的具体应用。
\end{statementbox}
